\section{Status supports and overlap scores}\label{sec:finite-partitions}

Use generic unary queries and common partitions whose cells are the full
fibers of linear coordinates. At orientation \(o\), let \(h_{o,v}(C)\) be the conditional
probability of status \(v\) in cell \(C\), and put
\(S_{o,v}=\{C:h_{o,v}(C)>0\}\).
If the query gives a cell zero mass, assign that cell the unconditioned
status law. This defines \(h_{o,v}\) everywhere. On the orientations
where we later extract predictions, every reference cell is attained.

The full status depends on at most \(m\) affine linear forms of the input,
where \(m\) is a construction constant. On any nonempty affine cell,
these forms are uniform on an affine space of dimension at most \(m\).
A possible status therefore has probability at least \(\theta=2^{-m}\).
The same argument applies to the unconditioned law on an empty cell. Conditioning on a diagonal bit
already recorded in the status can only increase its probability.
Thus
\[
 h_{o,v}(C)\ge\theta\,1_{S_{o,v}}(C).
\]
This bound needs no full-rank assumption on the coordinate map.

Write \(\mu(C)\) for the unary reference cell weight. Define
\begin{equation}\label{eq:table-overlap}
 O(v_A,v_B)=\sum_C\mu(C)1_{S_{o_A,v_A}}(C)1_{S_{o_B,v_B}}(C).
\end{equation}
For a complete numerical option \(\omega\), set
\begin{equation}\label{eq:finite-score}
 s_\omega(o_A,o_B)=1_{\{\text{flags pass}\}}
       \kappa_{A,\omega}\kappa_{B,\omega}
       \prod_i O_i(v_{A,i},v_{B,i}).
\end{equation}
The constants \(\kappa\) are the probabilities of the prescribed binary
patterns. There are finitely many options, and their positive \(\kappa\)'s
have a positive minimum. The numerical data defining an option can be enumerated before any
partition is chosen.

\subsection{From supports to accepted measures}

For an option with \(s\) positions put
\[
 f_o(k)=\kappa_\omega\theta^s
             \prod_i1_{S_{o,v_i}}(C_i(k_i)).
\]
At each position the corresponding indicator records whether the
required status is possible. The factor \(\theta^s\) accounts for its
minimum probability at all \(s\) positions. Multiply by the orientation
filter and average over a leaf to obtain \(f\).
Let \(d\) be the density of accepted keys before testing independence modulo \(F_0\).
The exact product expansion \eqref{eq:exact-product}, unary concentration,
and binary mixing give
\begin{equation}\label{eq:projected-error}
 \|(\Pi_Hf-\Pi_Hd)_+\|_1\le\epsilon
\end{equation}
in average over leaves. To see the direction of this inequality, test
against a nonnegative function on the common partition. Its product
expansion uses nonnegative unary cell tests. The support lower bound
applies to each factor. Binary mixing compares the resulting product
with the probability of accepting the tuple. Testing the set where the projected excess is
positive gives \eqref{eq:projected-error}.
The error can be made arbitrarily small, uniformly over the permitted
filters chosen using both leaves.

For functions of separate positions the Gram reference law satisfies
\[
 \int\prod_iF_i(k_i)\,d\nu
   =\prod_i\int F_i\,d\mu_i+O(2^{-cN}),\qquad |F_i|\le1.
\]
Expand the fixed cross-position Gram conditions and apply the Walsh bound.
It follows that \(\theta^{2s}\) times the expected score
\eqref{eq:finite-score} differs from \(\langle f,g\rangle_\nu\)
by an exponentially small error. The factor \(\theta^{2s}>0\) is fixed
before choosing the partition or the comparison accuracy.

% lax begin Lax793645.ProjectedOverlap
\begin{lemma}[Direct projection comparison]\label{lem:projected-mend}
On a finite probability space, let \(0\le f,g\le1\), \(d,e\ge0\),
and let \(\Pi\) be averaging on a partition. Then
\begin{equation}\label{eq:projected-mend}
 \langle f,g\rangle^2
 \le\langle\Pi d,\Pi e\rangle
       +\|(\Pi f-\Pi d)_+\|_1+\|(\Pi g-\Pi e)_+\|_1.
\end{equation}
\end{lemma}
% lax end
% lax begin Lax793645Proofs.Manuscript.ProjectedOverlap.one_sided_projected_mend
\begin{proof}
Put \(u=\Pi(fg)\). Since \(0\le u\le\Pi f,\Pi g\),
Cauchy--Schwarz gives
\(\langle f,g\rangle^2\le\E u^2\le\langle\Pi f,\Pi g\rangle\).
For \(a,b\in[0,1]\) and \(c,d\ge0\),
\(ab\le cd+(a-c)_++(b-d)_+\). Apply this pointwise and integrate.
\end{proof}
% lax end

The common-projection estimate from \Cref{sec:product-tests} bounds the
first term by actual overlap plus a small error, averaged over independent
leaves. Equation \eqref{eq:projected-error} bounds the remaining terms.
The acceptance filters are applied after drawing the two leaves.
Their dependence on the pair does not alter this averaging argument.

It remains to discard inputs whose base vectors are dependent modulo
the pins' protected base space \(F_0\). If \(d^G\le d\)
and \(\|d-d^G\|_1\le\delta_n\), then
\[
 \|(\Pi f-\Pi d^G)_+\|_1
 \le\|(\Pi f-\Pi d)_+\|_1+\delta_n.
\]
For the passage back to the product of the actual densities, use
\eqref{eq:bounded-overlap-transfer}. A bounded approximant for \(d\)
also approximates \(d^G\), with at most \(\delta_n\) additional
\(L^1\) error. By \eqref{eq:base-rejection} this error decreases
exponentially with \(n\).

% lax begin Lax793645.SupportScores
\begin{proposition}\label{prop:finite-score-bound}
Suppose all expected actual successful-query overlaps tend to zero along
a prepared sequence. For every fixed \(a>0\) and every \(\epsilon>0\),
common finite unary partitions can be chosen so that the probability of
a usable generic option with score at least \(a\) is at most \(\epsilon\),
eventually. Their sizes are bounded independently of \(n\).
They precede all permitted filters depending on the two leaves.
\end{proposition}
% lax end
% lax begin Lax793645Proofs.Manuscript.SupportScores.one_generic_lightmap
\begin{proof}
Let \(J_0\) be the number of options. For fixed leaves let
\(X_\omega=\langle f,g\rangle_\nu\). Equations \eqref{eq:projected-error}
and \eqref{eq:projected-mend}, followed by
\eqref{eq:bounded-overlap-transfer}, bound \(\E X_\omega^2\) by the
expected actual overlap plus an error that we can make arbitrarily
small. Since \(\E X_\omega\le(\E X_\omega^2)^{1/2}\), the product
estimate above makes \(\E s_\omega<a\epsilon/J_0\) for large \(n\).
Here we can choose the accuracy for all options at once because \(J_0\)
and the factors \(\theta^{2s}\) are fixed. Summing over the options
and applying \eqref{eq:score-markov} proves the claim.

\end{proof}
% lax end

% lax begin Lax793645.FiniteScores
For nonnegative scores \(s_\omega\) on any finite probability space,
\begin{equation}\label{eq:score-markov}
 a\,\PP\{\exists\omega:s_\omega\ge a\}
 \le\sum_\omega\E s_\omega\qquad(a>0).
\end{equation}
Indeed the indicator on the left is at most
\(a^{-1}\sum_\omega s_\omega\) pointwise.
% lax end

If supersaturation fails along a sequence, Proposition~\ref{prop:collision}
gives the hypothesis of Proposition~\ref{prop:finite-score-bound}.
We next show that the scalar equations force a positive score on a
positive fraction of pairs in the absence of common predictions.
